Model predictive control (MPC) for hexapods must optimize over high-dimensional hybrid dynamics while preserving physically valid joint branches. In practice, we observed that a warm-started Sequential Quadratic Programming (SQP) pipeline could converge to task-space-consistent yet branch-inconsistent tarsus configurations, after which runtime guards suppressed the commanded motion and destroyed forward progression. We addressed this failure by introducing a hard \emph{TarsusBranchConstraint} that encoded left/right tarsus sign admissibility directly in the optimal control problem, and by solving the resulting constrained MPC with an Interior-Point Method (IPM). Unlike the previous SQP path, the IPM backend kept the branch constraint active in the solver loop and excluded the forbidden branch from the feasible set rather than repairing it after optimization. In controlled short-horizon closed-loop experiments on an 18-actuator hexapod in Isaac Sim, the proposed method reduced branch-guard events from 275.9 to 0.0 per run, lowered tracking root-mean-square error (RMSE) from 0.242~m to 0.175~m, and maintained real-time operation with 4.82~ms average MPC solve time. Under the promoted endurance configuration, a repeated 20-second evaluation further achieved 0.306~m mean forward displacement over five runs with zero guard and reject events. Furthermore, in a 60-second endurance evaluation, the proposed method sustained 0.638~m forward displacement with 8.27$^\circ$ yaw drift, 0.011~m body-height range, and 4.63~ms average solve time. These results show that solver-visible physical branch constraints are critical for reliable multi-legged locomotion and that IPM offers a practical route to safe, real-time, and endurance-capable hexapod MPC.
